\textbf{6.}

Let $S$ be the collection of sets of the form
\[
(a,\infty)=\{x\in X:a<x\}
\]
and
\[
(-\infty,b)=\{x\in X:x<b\}.
\]

For two rays of the same type, we have
\[
(a,\infty)\cap(c,\infty)
=(\max\{a,c\},\infty)
\]
and
\[
(-\infty,b)\cap(-\infty,d)
=(-\infty,\min\{b,d\}).
\]

Also, for two rays of opposite types,
\[
(a,\infty)\cap(-\infty,b)=(a,b).
\]

Thus, finite intersections of elements of $S$ give
open intervals and open rays, which form a basis
for the order topology. Hence, $S$ is a subbasis
and generates the order topology on $X$.

\bigskip

\textbf{7.}

\textbf{(1)}

We prove this by induction on $k$.

For $k=1$, the result is clear since $A_1$ is
countable.

Suppose that
\[
A_1\times A_2\times\cdots\times A_k
\]
is countable.

Since $A_{k+1}$ is also countable and the product
of two countable sets is countable, we get
\[
(A_1\times\cdots\times A_k)\times A_{k+1}
\]
is countable.

Therefore, by induction,
\[
\prod_{i=1}^{k}A_i
\]
is countable for every $k\in\mathbb N$.

\medskip

\textbf{(2)}

No, the infinite Cartesian product is not
necessarily countable.

Consider
\[
A_n=\{0,1\},\qquad n\in\mathbb N.
\]

Suppose that
\[
\prod_{n=1}^{\infty}A_n
\]
is countable. Then we can list all its elements as
\[
x_1,x_2,x_3,\ldots
\]
where
\[
x_i=(x_{i1},x_{i2},\ldots).
\]

Define a new sequence $y=(y_1,y_2,\ldots)$ by
\[
y_n=1-x_{nn}.
\]

Clearly, $y$ belongs to the infinite product,
since every entry is either $0$ or $1$.

But $y$ is different from every $x_n$, since
their $n$-th entries are different.

This is a contradiction. Hence,
\[
\prod_{n=1}^{\infty}A_n
\]
is not necessarily countable.

\bigskip

\textbf{8.}

Suppose first that $\{x\}$ is open in the order
topology.

Then there exists a basic open set $B$ such that
\[
x\in B\subseteq\{x\}.
\]
Hence, $B=\{x\}$.

If $x$ is neither the minimum nor the maximum
element, then $B$ must be an open interval
$(a,b)$, where $a<x<b$.

Since $(a,b)=\{x\}$, there is no element between
$a$ and $x$, or between $x$ and $b$.
Thus, $a$ is the immediate predecessor and $b$
is the immediate successor of $x$.

If $x$ is the minimum element, then
\[
(-\infty,b)=\{x\},
\]
so $b$ is the immediate successor of $x$.

Similarly, if $x$ is the maximum element, then
\[
(a,\infty)=\{x\},
\]
so $a$ is the immediate predecessor of $x$.

Conversely, suppose that $x$ has an immediate
predecessor $a$ and an immediate successor $b$.
Then
\[
(a,b)=\{x\}.
\]
Since $(a,b)$ is open, $\{x\}$ is open.

If $x$ is the minimum element and has an
immediate successor $b$, then
\[
(-\infty,b)=\{x\},
\]
which is open.

Similarly, if $x$ is the maximum element and has
an immediate predecessor $a$, then
\[
(a,\infty)=\{x\},
\]
which is open.

Therefore, the given conditions are necessary
and sufficient for $\{x\}$ to be open.